% AUTO-GENERATED by scripts/figures/gen_cover.py — DO NOT EDIT BY HAND.
% Перегенерировать: python3 scripts/figures/gen_cover.py
% Векторная обложка книги (B5, full-bleed). Подключается в src/preamble.tex
% (после src/styles, где загружены tikz+fadings); рисуется \drawbookcover
% из \makebookcover (src/frontmatter/pre.tex). Опирается на уже загруженные
% tikz, tikzlibrary fadings, graphicx, eso-pic, paratype.
%
% Палитра намеренно холодная сине-стальная: чернильный синий +
% стальной синий; mesh-пятна — синяя пастель (родня открывателям частей).
\definecolor{CoverInk}{HTML}{0C2236}     % чернильный синий (тёмный конец)
\definecolor{CoverSteel}{HTML}{2F6E95}   % стальной синий (светлый конец)
\definecolor{CoverBloomA}{HTML}{CFE0F1}  % mesh: светлый стальной
\definecolor{CoverBloomB}{HTML}{DDE9F5}  % mesh: бледнее
\definecolor{CoverBloomC}{HTML}{BFD4E9}  % mesh: глубже
\tikzfading[name=coverbloom, inner color=transparent!8, outer color=transparent!100]

\newcommand{\coverplacetess}[3]{%  #1 cx  #2 cy  #3 scale (см)
  \coordinate (cv0) at ({#1+#3*(-0.3973)},{#2+#3*(-0.3141)});
  \coordinate (cv1) at ({#1+#3*(-1.0000)},{#2+#3*(-0.4712)});
  \coordinate (cv2) at ({#1+#3*(-0.0219)},{#2+#3*(-0.2432)});
  \coordinate (cv3) at ({#1+#3*(-0.3717)},{#2+#3*(-0.3572)});
  \coordinate (cv4) at ({#1+#3*(-0.3160)},{#2+#3*(0.3196)});
  \coordinate (cv5) at ({#1+#3*(-0.8416)},{#2+#3*(0.6381)});
  \coordinate (cv6) at ({#1+#3*(0.0580)},{#2+#3*(0.4259)});
  \coordinate (cv7) at ({#1+#3*(-0.2157)},{#2+#3*(0.8654)});
  \coordinate (cv8) at ({#1+#3*(0.1270)},{#2+#3*(-0.5097)});
  \coordinate (cv9) at ({#1+#3*(-0.1206)},{#2+#3*(-0.8862)});
  \coordinate (cv10) at ({#1+#3*(0.5841)},{#2+#3*(-0.4428)});
  \coordinate (cv11) at ({#1+#3*(0.7904)},{#2+#3*(-0.7994)});
  \coordinate (cv12) at ({#1+#3*(0.2086)},{#2+#3*(0.2005)});
  \coordinate (cv13) at ({#1+#3*(0.0433)},{#2+#3*(0.4809)});
  \coordinate (cv14) at ({#1+#3*(0.6625)},{#2+#3*(0.3122)});
  \coordinate (cv15) at ({#1+#3*(0.9415)},{#2+#3*(0.7443)});%
}

\newcommand{\coverdrawtess}{%
  \draw[draw=CoverGradHi!48!CoverGradLo,line cap=round,line join=round,line width=0.51pt,opacity=0.40] (cv9)--(cv13);
  \draw[draw=CoverGradHi!56!CoverGradLo,line cap=round,line join=round,line width=0.52pt,opacity=0.40] (cv12)--(cv13);
  \draw[draw=CoverGradHi!30!CoverGradLo,line cap=round,line join=round,line width=0.59pt,opacity=0.48] (cv5)--(cv13);
  \draw[draw=CoverGradHi!50!CoverGradLo,line cap=round,line join=round,line width=0.61pt,opacity=0.49] (cv8)--(cv9);
  \draw[draw=CoverGradHi!58!CoverGradLo,line cap=round,line join=round,line width=0.62pt,opacity=0.50] (cv8)--(cv12);
  \draw[draw=CoverGradHi!47!CoverGradLo,line cap=round,line join=round,line width=0.66pt,opacity=0.54] (cv4)--(cv12);
  \draw[draw=CoverGradHi!75!CoverGradLo,line cap=round,line join=round,line width=0.66pt,opacity=0.54] (cv13)--(cv15);
  \draw[draw=CoverGradHi!72!CoverGradLo,line cap=round,line join=round,line width=0.69pt,opacity=0.57] (cv12)--(cv14);
  \draw[draw=CoverGradHi!22!CoverGradLo,line cap=round,line join=round,line width=0.70pt,opacity=0.58] (cv1)--(cv9);
  \draw[draw=CoverGradHi!43!CoverGradLo,line cap=round,line join=round,line width=0.71pt,opacity=0.59] (cv0)--(cv8);
  \draw[draw=CoverGradHi!21!CoverGradLo,line cap=round,line join=round,line width=0.73pt,opacity=0.61] (cv4)--(cv5);
  \draw[draw=CoverGradHi!68!CoverGradLo,line cap=round,line join=round,line width=0.75pt,opacity=0.63] (cv8)--(cv10);
  \draw[draw=CoverGradHi!32!CoverGradLo,line cap=round,line join=round,line width=0.76pt,opacity=0.63] (cv0)--(cv4);
  \draw[draw=CoverGradHi!67!CoverGradLo,line cap=round,line join=round,line width=0.78pt,opacity=0.65] (cv9)--(cv11);
  \draw[draw=CoverGradHi!4!CoverGradLo,line cap=round,line join=round,line width=0.78pt,opacity=0.66] (cv1)--(cv5);
  \draw[draw=CoverGradHi!15!CoverGradLo,line cap=round,line join=round,line width=0.81pt,opacity=0.68] (cv0)--(cv1);
  \draw[draw=CoverGradHi!44!CoverGradLo,line cap=round,line join=round,line width=0.82pt,opacity=0.69] (cv4)--(cv6);
  \draw[draw=CoverGradHi!81!CoverGradLo,line cap=round,line join=round,line width=0.83pt,opacity=0.70] (cv10)--(cv14);
  \draw[draw=CoverGradHi!90!CoverGradLo,line cap=round,line join=round,line width=0.84pt,opacity=0.70] (cv14)--(cv15);
  \draw[draw=CoverGradHi!68!CoverGradLo,line cap=round,line join=round,line width=0.86pt,opacity=0.72] (cv6)--(cv14);
  \draw[draw=CoverGradHi!40!CoverGradLo,line cap=round,line join=round,line width=0.88pt,opacity=0.74] (cv0)--(cv2);
  \draw[draw=CoverGradHi!24!CoverGradLo,line cap=round,line join=round,line width=0.90pt,opacity=0.76] (cv5)--(cv7);
  \draw[draw=CoverGradHi!64!CoverGradLo,line cap=round,line join=round,line width=0.91pt,opacity=0.78] (cv2)--(cv10);
  \draw[draw=CoverGradHi!84!CoverGradLo,line cap=round,line join=round,line width=0.93pt,opacity=0.79] (cv10)--(cv11);
  \draw[draw=CoverGradHi!93!CoverGradLo,line cap=round,line join=round,line width=0.93pt,opacity=0.80] (cv11)--(cv15);
  \draw[draw=CoverGradHi!51!CoverGradLo,line cap=round,line join=round,line width=0.94pt,opacity=0.80] (cv2)--(cv6);
  \draw[draw=CoverGradHi!68!CoverGradLo,line cap=round,line join=round,line width=0.97pt,opacity=0.83] (cv7)--(cv15);
  \draw[draw=CoverGradHi!46!CoverGradLo,line cap=round,line join=round,line width=0.99pt,opacity=0.85] (cv6)--(cv7);
  \draw[draw=CoverGradHi!16!CoverGradLo,line cap=round,line join=round,line width=0.99pt,opacity=0.85] (cv1)--(cv3);
  \draw[draw=CoverGradHi!40!CoverGradLo,line cap=round,line join=round,line width=1.06pt,opacity=0.91] (cv2)--(cv3);
  \draw[draw=CoverGradHi!60!CoverGradLo,line cap=round,line join=round,line width=1.07pt,opacity=0.93] (cv3)--(cv11);
  \draw[draw=CoverGradHi!35!CoverGradLo,line cap=round,line join=round,line width=1.10pt,opacity=0.96] (cv3)--(cv7);
  \fill[CoverGradHi!30!CoverGradLo,opacity=0.65] (cv0) circle (0.040);
  \fill[CoverGradHi!0!CoverGradLo,opacity=0.70] (cv1) circle (0.042);
  \fill[CoverGradHi!49!CoverGradLo,opacity=0.83] (cv2) circle (0.046);
  \fill[CoverGradHi!31!CoverGradLo,opacity=1.00] (cv3) circle (0.052);
  \fill[CoverGradHi!34!CoverGradLo,opacity=0.61] (cv4) circle (0.039);
  \fill[CoverGradHi!8!CoverGradLo,opacity=0.61] (cv5) circle (0.039);
  \fill[CoverGradHi!53!CoverGradLo,opacity=0.78] (cv6) circle (0.045);
  \fill[CoverGradHi!39!CoverGradLo,opacity=0.91] (cv7) circle (0.049);
  \fill[CoverGradHi!56!CoverGradLo,opacity=0.52] (cv8) circle (0.036);
  \fill[CoverGradHi!44!CoverGradLo,opacity=0.45] (cv9) circle (0.034);
  \fill[CoverGradHi!79!CoverGradLo,opacity=0.73] (cv10) circle (0.043);
  \fill[CoverGradHi!90!CoverGradLo,opacity=0.85] (cv11) circle (0.047);
  \fill[CoverGradHi!60!CoverGradLo,opacity=0.47] (cv12) circle (0.034);
  \fill[CoverGradHi!52!CoverGradLo,opacity=0.34] (cv13) circle (0.030);
  \fill[CoverGradHi!83!CoverGradLo,opacity=0.67] (cv14) circle (0.041);
  \fill[CoverGradHi!97!CoverGradLo,opacity=0.74] (cv15) circle (0.043);%
}

\newcommand{\drawbookcover}{%
  \begin{tikzpicture}
    \useasboundingbox (0,0) rectangle (17.6,25);
    \clip (0,0) rectangle (17.6,25);
    \fill[white] (0,0) rectangle (17.6,25);
    % чертёжная сетка (blueprint): тонкая 1 см + усиленная 5 см, чернилами
    \draw[CoverInk,line width=0.2pt,opacity=0.05,step=1cm] (0,0) grid (17.6,25);
    \draw[CoverInk,line width=0.4pt,opacity=0.09,step=5cm] (0,0) grid (17.6,25);
    % мягкие синие mesh-пятна по углам
    \fill[CoverBloomA,path fading=coverbloom] (15.2,22.0) circle (8.0);
    \fill[CoverBloomB,path fading=coverbloom] (1.8,3.5) circle (8.5);
    \fill[CoverBloomC,path fading=coverbloom] (16.6,10.5) circle (5.5);
    % гиперкуб: стальной (слева) → чернильный (справа)
    \colorlet{CoverGradLo}{CoverSteel}\colorlet{CoverGradHi}{CoverInk}%
    \coverplacetess{8.8}{16.7}{4.2}
    \coverdrawtess
    % угловые реперные метки
    \begin{scope}[draw=CoverInk!55,line width=0.5pt,opacity=0.85]
      \draw (0.85,0.85)--+(0.5,0); \draw (0.85,0.85)--+(0,0.5);
      \draw (16.75,0.85)--+(-0.5,0); \draw (16.75,0.85)--+(0,0.5);
      \draw (0.85,24.15)--+(0.5,0); \draw (0.85,24.15)--+(0,-0.5);
      \draw (16.75,24.15)--+(-0.5,0); \draw (16.75,24.15)--+(0,-0.5);
    \end{scope}
    % верхняя метка
    \node[anchor=center,text=CoverInk!70,font=\sffamily] at (8.8,23.7) {\fontsize{9}{9}\selectfont O\kern0.22em P\kern0.22em E\kern0.22em N\kern0.22em \kern0.22em S\kern0.22em O\kern0.22em U\kern0.22em R\kern0.22em C\kern0.22em E};
    % титул во всю ширину полосы (сплошной чернильный синий)
    \node[anchor=center] at (8.8,7.3) {\resizebox{15.8cm}{!}{\sffamily\bfseries\color{CoverInk}ИНЖЕНЕРИЯ ПЛАТЕЖЕЙ}};
    % подзаголовок
    \node[anchor=center,text=CoverInk!82,font=\sffamily,align=center,text width=15cm] at (8.8,5.6) {\fontsize{11}{15}\selectfont Открытая книга о карточных платежах, СБП и платёжной инфраструктуре};
    % год
    \node[anchor=center,text=CoverInk!70,font=\sffamily\bfseries] at (8.8,1.45) {\fontsize{11}{11}\selectfont 2026};
  \end{tikzpicture}%
}
